\documentclass[12pt]{article}
\usepackage[utf8]{inputenc}
\usepackage[T1]{fontenc}
\usepackage{hyperref}
\title{Aghwan Language Resources for K-12 Curriculum Development}
\author{Compiled by Claude Sonnet 5.5 (claude-sonnet-5-5) for Andr\'as Kornai}
\date{}
\begin{document}
\maketitle

\section{Introduction}
Aghwan (ISO 639-3 xag; Glottocode aghw1237) is the extinct language usually called Caucasian Albanian, also Old Udi or Aluan. It is a Northeast Caucasian (Nakh-Daghestanian) language. Glottolog places it in the Udi-Aghwan subgroup of the Lezgic branch, and it is treated as the predecessor of the endangered Udi language. It was spoken in Caucasian Albania, from southern Dagestan to Azerbaijan, and had its own alphabet of about 52 letters. The main surviving material is the palimpsests at Mt.\ Sinai and inscriptions from Mingachevir, Azerbaijan. Some of the Lezgic-origin claims are suppositions and not established fact.

\section{Primary Educational Resources}
No school primers, textbooks or programmes for Aghwan were found. The language is extinct, and the living relative Udi is the more likely basis for classroom work, though the search did not turn up Udi teaching materials.

\section{Linguistic and Reference Materials}
\begin{itemize}
\item Wikipedia on the Caucasian Albanian language: \url{https://en.wikipedia.org/wiki/Caucasian_Albanian_language}
\item Glottolog entry for Udi, the descendant language: \url{https://glottolog.org/resource/languoid/id/udia1235}
\item Grambank entry for Aghwan: \url{https://grambank.clld.org/languages/aghw1237}
\item Omniglot page on Caucasian Albanian numbers: \url{https://www.omniglot.com/language/numbers/caucasianalbanian.htm}
\end{itemize}

\section{Religious and Community Resources}
The surviving texts are mostly Christian, in the Sinai palimpsests. No community organisation or modern religious edition with a URL was found in the search results.

\section{Digital and Audio Resources}
No audio, corpora or software were found.

\section{Curriculum Development Context}
There are no speakers, and the corpus is small. A curriculum could cover the script, the history of Caucasian Albania and the relation to Udi, as a heritage and history unit. Detailed language content would need specialist sources beyond what the search returned.

\end{document}